% Generated from proposal.md; edit the Markdown, then rebuild.
% Markdown SHA256: e2586193fe5e5a15b9d3d1a72684eff3cf94f74b4730a72e62df6e7459ab498c
% Shared template for Pandoc -> XeLaTeX/Tectonic. Content lives in the .md files.
\documentclass[11pt,a4paper]{article}
\usepackage[margin=22mm,headheight=14pt]{geometry}
\usepackage{fontspec}
\setmainfont{texgyretermes}[Extension=.otf,UprightFont=*-regular,BoldFont=*-bold,ItalicFont=*-italic,BoldItalicFont=*-bolditalic]
\setsansfont{texgyreheros}[Extension=.otf,UprightFont=*-regular,BoldFont=*-bold,ItalicFont=*-italic,BoldItalicFont=*-bolditalic]
\setmonofont{texgyrecursor}[Extension=.otf,UprightFont=*-regular,BoldFont=*-bold,ItalicFont=*-italic,BoldItalicFont=*-bolditalic,Scale=MatchLowercase]
\usepackage{amsmath,amssymb}
\usepackage{array,booktabs,longtable,calc}
\usepackage{graphicx}
\usepackage{xurl}
\usepackage[unicode,hidelinks]{hyperref}
\usepackage{bookmark}
\usepackage{enumitem}
\usepackage{titlesec}
\usepackage{fancyhdr}
\usepackage{microtype}
\hypersetup{pdftitle={Đề cương tuần
1},pdfauthor={Cao Văn Hà; Hà Mạnh Thế; Trần Duy Hưng},pdflang={vi-VN}}
\urlstyle{tt}
\setlength{\parindent}{0pt}
\setlength{\parskip}{5pt plus 1pt minus 1pt}
\linespread{1.08}
\setlist{nosep,leftmargin=*,topsep=4pt}
\setcounter{secnumdepth}{0}
\titleformat{\section}{\large\bfseries}{\thesection}{0pt}{}
\titleformat{\subsection}{\normalsize\bfseries}{\thesubsection}{0pt}{}
\titlespacing*{\section}{0pt}{12pt}{5pt}
\titlespacing*{\subsection}{0pt}{9pt}{4pt}
\setlength{\tabcolsep}{4pt}
\renewcommand{\arraystretch}{1.16}
\setlength{\LTpre}{6pt}
\setlength{\LTpost}{6pt}
\setlength{\emergencystretch}{2em}
\providecommand{\tightlist}{\setlength{\itemsep}{0pt}\setlength{\parskip}{0pt}}
\providecommand{\pandocbounded}[1]{#1}
\pagestyle{fancy}
\fancyhf{}
\fancyhead[L]{\small Đề tài 9 — Hồ sơ tuần 1}
\fancyhead[R]{\small Bản nháp}
\fancyfoot[C]{\thepage}
\renewcommand{\headrulewidth}{0.3pt}
\begin{document}
{\LARGE\bfseries Đề cương tuần 1\par}
\vspace{4pt}
{\large Kiến trúc Transformer và tinh chỉnh mô hình ngôn ngữ\par}
\vspace{4pt}
\textbf{Ứng dụng:} Hỏi đáp thủ tục hộ tịch–chứng thực tại Hà Nội.\\
\textbf{Nhóm:} Cao Văn Hà (trưởng nhóm), Hà Mạnh Thế, Trần Duy Hưng.\\
\textbf{Phiên bản:} Bản nháp tuần 1 — có audit nghiên cứu, chưa có dữ liệu huấn luyện được phê duyệt.
\vspace{5pt}\hrule\vspace{8pt}
\begin{quote}
\textbf{Trạng thái:} bản nháp tuần 1, đã thống nhất phạm vi dữ liệu và
phương án khảo sát/chấm. Những nội dung còn ghi ``đề xuất'' chưa phải dữ
liệu đã thu thập, thí nghiệm đã chạy hay quyết định cuối cùng. Mục tiêu
độ dài nội bộ ban đầu là khoảng 5--7 trang đề cương và 3--5 trang kế
hoạch dữ liệu; không phải yêu cầu của giảng viên. Bản PDF giữ đầy đủ nội
dung Markdown, kể cả bảng và schema, nên có thể dài hơn mục tiêu này.

\textbf{Căn cứ đề bài:} Đề tài 9 yêu cầu cài đặt Transformer, LoRA và
hàm mất mát DPO; huấn luyện mô hình ký tự nhỏ; sau đó xây dựng trợ lý
hỏi đáp thủ tục hành chính với tối thiểu 2.000 cặp chỉ dẫn có trích dẫn,
ít nhất 500 cặp ưu tiên và đánh giá mô hình gốc/SFT/SFT+DPO. Phần D có
tỷ trọng 30/100; đề bài cũng yêu cầu kết quả trung bình ± độ lệch chuẩn
trên ít nhất 5 lần chạy hoặc 5 phần gập, và quản lý phiên bản nguồn thủ
tục.
\end{quote}

\section{1. Mục tiêu, phạm vi và người
dùng}\label{mux1ee5c-tiuxeau-phux1ea1m-vi-vuxe0-ngux1b0ux1eddi-duxf9ng}

\subsection{1.1. Vấn đề và mục
tiêu}\label{vux1ea5n-ux111ux1ec1-vuxe0-mux1ee5c-tiuxeau}

Nhóm xây dựng và đánh giá một \textbf{trợ lý hỏi đáp bằng tiếng Việt về
thủ tục hành chính công trong phạm vi được chốt}, nhằm trả lời có căn cứ
cho các câu hỏi thường gặp như thành phần hồ sơ, trình tự, thời hạn, lệ
phí và căn cứ pháp lý. Câu trả lời được kỳ vọng nêu rõ trích dẫn tới
đoạn nguồn đã lưu; khi câu hỏi nằm ngoài phạm vi dữ liệu hoặc không đủ
căn cứ, hệ thống phải từ chối hoặc nêu giới hạn thay vì tự suy đoán.

Mục tiêu học thuật là nối liền bốn phần bắt buộc của Đề tài 9:

\begin{itemize}
\tightlist
\item
  \textbf{A --- phát triển phương pháp:} tự cài đặt khối Transformer
  giải mã và các thành phần được yêu cầu; tự cài đặt LoRA và mất mát
  DPO.
\item
  \textbf{B --- lý thuyết:} trình bày các chứng minh về attention, mã
  hóa vị trí, LoRA, DPO và cross-entropy/PPL theo đề bài.
\item
  \textbf{C --- bài tập và kiểm chứng:} giải các bài tính attention, độ
  phức tạp, LoRA/QLoRA và giá trị--gradient DPO; kiểm chứng số học các
  kết quả liên quan.
\item
  \textbf{D --- ứng dụng:} tạo bộ dữ liệu có nguồn gốc rõ ràng, SFT một
  mô hình nguồn mở phù hợp tiếng Việt và tiếp tục tối ưu DPO từ dữ liệu
  ưu tiên do người đánh giá lập. Phần B sẽ liệt kê và chứng minh đầy đủ
  sáu mệnh đề của đề bài: (B.1) phương sai \(q^\top k\) và hệ số
  \(1/\sqrt{d_k}\); (B.2) đẳng biến hoán vị khi không có mã hóa vị trí;
  (B.3) ma trận quay khối cho mã hóa hình sin; (B.4) bốn tính chất LoRA;
  (B.5) nghiệm tối ưu DPO có chính quy KL từ Bradley--Terry; và (B.6)
  tương đương giữa tối thiểu cross-entropy kỳ vọng và KL, cùng liên hệ
  PPL. Phần A xác định rõ causal mask gán logit token tương lai giá trị
  \(-\infty\) trước softmax. Khi chứng minh B.5, báo cáo phải nêu các
  giả thiết như \(\beta>0\), hỗ trợ của \(\pi_{ref}\) bao phủ các đáp án
  đang xét và hằng số chuẩn hóa \(Z(x)\) hữu hạn/dương; không dùng công
  thức ngoài các điều kiện đó. Bài C.4 sẽ tính giá trị \textbf{và
  gradient} DPO từ log-xác suất chính sách/tham chiếu cho
  \nolinkurl{chosen} và \nolinkurl{rejected}, rồi phân tích ảnh hưởng
  của \(\beta\).
\end{itemize}

Với B.2, chứng minh trước cho attention không mask; nếu có mask \(M\),
phải đồng thời biến đổi thành \(PMP^\top\). Giữ nguyên causal mask nói
chung không bảo toàn đẳng biến với mọi hoán vị. Mọi chứng minh phải nêu
kích thước tensor, tính độc lập cần thiết và dùng log tự nhiên.

Mục tiêu đánh giá nội bộ có thể hướng tới mức 85--90/100, nhưng đây
\textbf{không phải cam kết điểm số}. Điểm chỉ có thể được xác định bởi
hội đồng sau khi xem toàn bộ chứng minh, mã nguồn, dữ liệu, tính tái lập
và bảo vệ.

\subsection{1.2. Phạm vi có chủ ý
hẹp}\label{phux1ea1m-vi-cuxf3-chux1ee7-uxfd-hux1eb9p}

Người dùng đã đồng ý chọn \textbf{Hà Nội}, tập trung \textbf{hộ
tịch--chứng thực}, hướng tới sinh viên, người trẻ và gia đình. Vòng
nghiên cứu đầu đã kiểm 9 nguồn cho 5 họ thủ tục ứng viên và chuẩn bị 20
QA nháp có dẫn đoạn nguồn. Chưa QA nào được con người phê duyệt hoặc đủ
điều kiện huấn luyện; chi tiết nằm trong sổ kiểm nguồn và báo cáo pilot.
Mục tiêu 5--10 thủ tục/20--50 QA là mốc thăm dò nội bộ, không phải
ngưỡng bắt buộc của giảng viên.

Trong phạm vi đã chọn, sản phẩm tập trung vào hỏi đáp một lượt có trích
dẫn, dựa trên mô tả thủ tục và văn bản/cổng công khai được phép sử dụng.
Sản phẩm không thay thế tư vấn pháp lý, không đưa ra quyết định hành
chính, không xử lý hồ sơ thật và không yêu cầu người dùng cung cấp dữ
liệu cá nhân.

Ngoài phạm vi của đề cương này:

\begin{itemize}
\tightlist
\item
  Không triển khai RAG; RAG chỉ được nhắc như hướng mở rộng mà đề bài đã
  gợi ý cho Đề tài 11.
\item
  Không triển khai PPO hoặc bất kỳ vòng học tăng cường trực tuyến nào;
  nhánh ưu tiên dùng DPO ngoại tuyến theo đề bài.
\item
  Không khẳng định cập nhật tri thức thời gian thực. Mỗi câu trả lời
  phải mang theo căn cứ/phiên bản nguồn đã biết; thủ tục thay đổi cần
  thu thập lại và gán phiên bản mới.
\end{itemize}

\subsection{1.3. Người dùng và tình huống sử
dụng}\label{ngux1b0ux1eddi-duxf9ng-vuxe0-tuxecnh-huux1ed1ng-sux1eed-dux1ee5ng}

Người dùng mục tiêu là người dân hoặc sinh viên cần \textbf{định hướng
ban đầu} về một thủ tục trong phạm vi được công bố. Họ có thể hỏi ``cần
những giấy tờ gì?'', ``nộp theo các bước nào?'', ``thời hạn/chi phí là
bao lâu?'' hoặc ``căn cứ nào nói điều này?''. Trợ lý phải ưu tiên câu
trả lời ngắn, dễ kiểm tra và liên kết/trích đoạn nguồn hơn là giọng điệu
tự tin.

Người dùng không nên dùng hệ thống cho tình huống khẩn cấp, tranh chấp,
trường hợp cá nhân có ngoại lệ hoặc thủ tục ngoài nhóm đã công bố. Với
các trường hợp đó, mẫu trả lời cần chỉ ra rằng người dùng phải kiểm tra
cổng/cơ quan có thẩm quyền. Việc khảo sát người dùng thật để lập tập
test chỉ diễn ra khi có văn bản đồng ý, có thể rút lại theo quy trình
khảo sát, và dữ liệu được ẩn danh trước khi lưu.

\section{2. Kiến trúc công việc: hai nhánh và cầu nối
A--D}\label{kiux1ebfn-truxfac-cuxf4ng-viux1ec7c-hai-nhuxe1nh-vuxe0-cux1ea7u-nux1ed1i-ad}

\subsection{2.1. Nhánh 1 --- nền tảng tự cài đặt cho Phần
A--C}\label{nhuxe1nh-1-nux1ec1n-tux1ea3ng-tux1ef1-cuxe0i-ux111ux1eb7t-cho-phux1ea7n-ac}

Nhánh này phục vụ yêu cầu ``từ đầu'', không dùng mô-đun attention dựng
sẵn cho phần cài đặt cốt lõi. Các thành phần dự kiến gồm: scaled
dot-product attention, multi-head attention, causal mask, positional
encoding hình sin và một block Transformer decoder bằng PyTorch; mô hình
ngôn ngữ mức ký tự nhỏ; lớp LoRA bao quanh linear layer với cơ chế hợp
nhất trọng số khi suy luận; và hàm mất mát DPO.

Mô hình ký tự nhỏ chỉ dùng văn bản tiếng Việt thuộc phạm vi công cộng
hoặc giấy phép mở đã được kiểm tra. Nhóm sẽ tách train/validation riêng
và báo cáo PPL validation. Đây là bài kiểm chứng kiến trúc/phương pháp,
không phải mô hình trợ lý Phần D và không được dùng văn bản hành chính
chưa có quyền sử dụng.

Đã audit kỹ thuật bản Truyện Kiều Wikisource revision \nolinkurl{206607}
trong kho cục bộ riêng: trích 3.254 dòng, chia tập liên tục, học từ vựng
từ train và kiểm cửa sổ không vượt ranh giới. Đây là kiểm tra dữ liệu
nghiên cứu, \textbf{không phải huấn luyện hoặc phê duyệt quyền dùng}.
Trước huấn luyện phải hoàn tất review quyền của bản số hóa; không đưa
corpus vào Git khi chưa được phép.

Cần kiểm chứng từng hạng mục có thể quan sát được: causal mask không lộ
token tương lai; shape/head của attention đúng; tham số LoRA trainable
bằng \(r(d+k)\); hợp nhất LoRA bảo toàn đầu ra trong sai số số học; hàm
DPO có dấu và gradient nhất quán với công thức. Các bài C.1--C.4, gồm
C.4 về giá trị--gradient DPO và ảnh hưởng của \(\beta\), sẽ được giải và
dùng làm phép kiểm tra độc lập, không thay cho đánh giá ứng dụng.

\subsection{2.2. Nhánh 2 --- ứng dụng SFT rồi DPO cho Phần
D}\label{nhuxe1nh-2-ux1ee9ng-dux1ee5ng-sft-rux1ed3i-dpo-cho-phux1ea7n-d}

Nhánh ứng dụng bắt đầu từ dữ liệu chỉ dẫn có trích dẫn và một mô hình
nền mã nguồn mở hỗ trợ tiếng Việt. Quy trình dự kiến: khóa phiên bản
nguồn và chia dữ liệu; SFT; sinh hai ứng viên trả lời từ checkpoint SFT
cho mỗi prompt; hai người đánh giá độc lập chọn câu trả lời tốt hơn theo
rubric; phân xử bất đồng; sau đó DPO với checkpoint SFT đóng băng làm mô
hình tham chiếu.

\textbf{Qwen2.5-1.5B-Instruct} là một ứng viên, không phải lựa chọn cuối
cùng. Lựa chọn mô hình sẽ phụ thuộc giấy phép dùng cho nghiên cứu, khả
năng tiếng Việt, tokenizer, khả năng nạp với QLoRA và kết quả pilot.
Người dùng đã cung cấp kết quả tự kiểm tra trong một phiên Colab trước:
T4 15.360 MiB, \nolinkurl{torch 2.14.1+cu130}, CUDA của PyTorch
\nolinkurl{13.0} và nhân ma trận CUDA thành công; agent chưa thực hiện
phép đo GPU này và phiên mới có thể khác. Không coi đó là bảo đảm đủ
VRAM: cần pilot SFT và pilot DPO riêng, với cách quản lý mô hình tham
chiếu tương thích bộ nhớ. Checkpoint, manifest và trạng thái resume dự
kiến lưu Drive nếu có quyền và đủ dung lượng; cần kiểm tra resume thực
tế, không cam kết không gián đoạn.

SFT pilot phải xuất hiện sớm cuối tuần 1 hoặc đầu tuần 2 để có
checkpoint tạo ứng viên cho dữ liệu DPO. Không được giả định đã có 500
cặp ưu tiên trước khi hai người chấm hoàn tất quy trình. Mục tiêu bảo
thủ để khóa kế hoạch là đủ \textbf{2.000 QA train} và \textbf{500
preference pairs train}, cộng validation riêng; con số thực chỉ được báo
sau khi thu thập, kiểm tra và phân chia.

\subsection{2.3. Cầu nối bắt buộc từ A sang
D}\label{cux1ea7u-nux1ed1i-bux1eaft-buux1ed9c-tux1eeb-a-sang-d}

Cầu nối không dừng ở việc đặt hai nhánh cạnh nhau. Trên
\textbf{minibatch thật của dữ liệu Phần D}, nhóm sẽ dùng cùng prompt,
\nolinkurl{chosen}, \nolinkurl{rejected}, attention mask và token mask
trả lời để tính log-xác suất chuỗi dưới mô hình chính sách và mô hình
tham chiếu. Với vanilla DPO, log-xác suất mỗi đáp án là \textbf{tổng}
log-prob các token thuộc phần trả lời đã mask, không phải trung bình
theo token như phép tính PPL; loss là trung bình các loss ví dụ theo
batch. Hàm tự cài đặt phải tính:

\[L_{DPO}=-\frac{1}{B}\sum_i\log\sigma\left(\beta\left[(\log\pi_\theta(y_w|x)-\log\pi_{ref}(y_w|x))-(\log\pi_\theta(y_l|x)-\log\pi_{ref}(y_l|x))\right]\right).\]

Đối chiếu với thư viện DPO được phép dùng ở Phần D sẽ ghi lại: tổng
log-prob \nolinkurl{chosen}/\nolinkurl{rejected} cho từng ví dụ, reward
margin/loss theo minibatch, masking, reduction và gradient. Với cùng
minibatch và trạng thái tham số ban đầu, đối chiếu gradient autograd
theo logit có mask hoặc tham số adapter có thể so sánh được; không chỉ
so scalar loss. Đầu vào phải là log-prob thực từ forward pass,
\textbf{không} là số ngẫu nhiên hay ví dụ làm sẵn. Sai khác chỉ được kết
luận sau khi cố định tokenizer, padding, truncation, prompt template,
dấu margin, token mask, cách tổng token trả lời và batch reduction. Kết
quả đối chiếu là bằng chứng rằng thành phần A tham gia trực tiếp vào
pipeline D, như yêu cầu chung của đề bài.

Trong pipeline D, LoRA tự cài đặt có thể được dùng làm phép kiểm tra
trên linear toy/mini-model; việc fine-tune quy mô Phần D có thể dùng thư
viện chuẩn theo quy định Phần D. Nếu hai cách không thể so sánh cùng
kiến trúc/tokenizer, báo cáo phải nói rõ giới hạn thay vì suy diễn tương
đương.

\section{3. Dữ liệu, chất lượng và đạo
đức}\label{dux1eef-liux1ec7u-chux1ea5t-lux1b0ux1ee3ng-vuxe0-ux111ux1ea1o-ux111ux1ee9c}

Kế hoạch vận hành chi tiết nằm trong \href{data-plan.pdf}{data-plan.md}.
Tóm tắt các yêu cầu không được lược bỏ:

\begin{enumerate}
\def\labelenumi{\arabic{enumi}.}
\tightlist
\item
  \textbf{Nguồn thủ tục:} chỉ thu thập trang/tệp công khai mà việc sử
  dụng được phép, tuân thủ điều khoản, robots.txt và giới hạn tần suất.
  Với mỗi đoạn, lưu URL/cơ quan, thời điểm truy xuất, phiên bản/ngày
  hiệu lực khi có, giấy phép/điều khoản đã kiểm tra và span nguồn có thể
  trích lại. Không suy diễn nguồn hành chính trước khi địa phương và
  nhóm thủ tục được chốt.
\item
  \textbf{SFT:} đề bài yêu cầu ít nhất 2.000 cặp QA có citation và kiểm
  định thủ công ít nhất 20\%. Citation phải trỏ tới span nguồn được lưu,
  không chỉ tới tên website. Train, validation và test tách theo nhóm
  thủ tục/nguồn để giảm rò rỉ gần trùng lặp.
\item
  \textbf{DPO:} đề bài yêu cầu tối thiểu 500 cặp so sánh, mỗi cặp là hai
  câu trả lời SFT cho cùng prompt do ít nhất hai người đánh giá độc lập.
  Rubric gồm đúng nội dung, đầy đủ, trích dẫn đúng và từ chối đúng phạm
  vi. Cặp bất đồng được phân xử sau khi đã lưu nhãn độc lập; báo cáo tỷ
  lệ đồng thuận và cách phân xử.
\item
  \textbf{Test người dùng:} câu hỏi khảo sát thật phải tách hoàn toàn
  với train/validation, thu thập có consent bằng văn bản và ẩn danh
  trước khi lưu. Nếu chưa có quyền tiếp cận người khảo sát, không được
  gọi tập test là ``câu hỏi thật''; phải ghi rõ là thiếu điều kiện.
\end{enumerate}

Các nguồn, dữ liệu, phiên bản và hướng dẫn gán nhãn cần được đóng gói
thành datasheet; mô hình cuối cùng cần model card. Cả hai là sản phẩm
nộp được yêu cầu, không chỉ là ghi chú nội bộ.

\section{4. Huấn luyện, chọn siêu tham số và đánh
giá}\label{huux1ea5n-luyux1ec7n-chux1ecdn-siuxeau-tham-sux1ed1-vuxe0-ux111uxe1nh-giuxe1}

\subsection{4.1. Lựa chọn và thí nghiệm có kiểm
soát}\label{lux1ef1a-chux1ecdn-vuxe0-thuxed-nghiux1ec7m-cuxf3-kiux1ec3m-souxe1t}

Mô hình nền chỉ được chốt sau khi kiểm tra giấy phép và pilot SFT. Các
cấu hình LoRA/QLoRA, hạng \(r\), \(\alpha\), learning rate, số epoch, độ
dài chuỗi, batch/gradient accumulation, \(\beta\) DPO và giải mã sẽ được
ghi bằng manifest. Không khóa giá trị số chỉ vì chúng phổ biến: siêu
tham số chỉ được chọn trên validation theo quy định đề bài; test chỉ
dùng cho báo cáo cuối.

DPO dùng checkpoint SFT đã chọn làm \nolinkurl{ref}, đóng băng trong một
run. Với seed \(s\), thiết kế run ghép là \nolinkurl{SFT_s → DPO_s},
trong đó \nolinkurl{ref = SFT_s}; preference dataset, split, rubric và
candidate đã phân xử được khóa chung từ pilot trước khi chạy 5 seed, nên
không đòi hỏi gán nhãn lại 500 cặp cho mỗi seed. Nhóm cần lưu
fingerprint checkpoint, tokenizer, prompt template và split dữ liệu để
các run sau tái lập đúng cặp so sánh. Dữ liệu validation DPO tách riêng
khỏi 500 cặp train đề xuất; không tối ưu \(\beta\) trên test.

Nếu ngân sách GPU không đủ cho 5 run độc lập của cấu hình cuối, nhóm sẽ
xin điều chỉnh \textbf{với giảng viên} và ghi rõ tác động; thỏa thuận
nội bộ của nhóm không thay thế điều chỉnh đó. \textbf{Không} âm thầm
giảm số run. Năm lần sinh câu trả lời từ cùng checkpoint không phải năm
run; chúng chỉ đo biến thiên decoding nếu thiết kế có chủ đích.

\subsection{4.2. So sánh và chỉ
số}\label{so-suxe1nh-vuxe0-chux1ec9-sux1ed1}

So sánh tối thiểu là mô hình gốc, SFT và SFT+DPO, trên cùng tập test
khóa và prompt template tương ứng. Đánh giá câu trả lời gồm:

\begin{itemize}
\tightlist
\item
  độ chính xác thông tin khi đối chiếu span nguồn;
\item
  tính đầy đủ với yêu cầu trong câu hỏi;
\item
  tỷ lệ citation đúng (citation có thật và hỗ trợ đúng mệnh đề);
\item
  tỷ lệ ảo giác/khẳng định không có căn cứ;
\item
  tỷ lệ từ chối đúng cho câu hỏi ngoài phạm vi hoặc thiếu căn cứ;
\item
  đánh giá mù theo cặp của người chấm, tỷ lệ thắng và khoảng tin cậy.
\end{itemize}

Các run SFT/DPO độc lập dùng 5 seed ghép như trên; đơn vị tổng hợp là
run, báo cáo mean và \textbf{sample standard deviation} (\(n-1\)), không
nhân bản mỗi prompt hay mỗi generation thành run. Baseline mô hình gốc
được khóa theo một checkpoint/config cố định và không được gán sample SD
bằng cách lặp suy luận/nhân bản output; nếu có ngẫu nhiên decoding phải
báo giao thức riêng, không gọi đó là biến thiên huấn luyện. Khoảng tin
cậy của đánh giá mù phải bootstrap theo \textbf{đơn vị câu hỏi}; nếu
khảo sát có nhiều câu từ cùng người, bootstrap theo \textbf{cụm người
khảo sát}. Không xem các quan sát seed × prompt là độc lập để làm CI hẹp
giả tạo. Số câu có chấm, số cụm, cách bootstrap và seed phải được báo
cáo.

Không có số liệu nào trong đề cương này là kết quả. Chỉ sau khi test
được khóa, chấm mù và lưu biểu mẫu mới được báo cáo chỉ số/CI.

\section{5. Cách ba thành viên phối
hợp}\label{cuxe1ch-ba-thuxe0nh-viuxean-phux1ed1i-hux1ee3p}

Nhóm gồm \textbf{Cao Văn Hà, Hà Mạnh Thế và Trần Duy Hưng}. Hà là trưởng
nhóm, thực hiện chính xuyên suốt A--D, tổng hợp và nộp bài. Không giao
cố định từng phần cho từng thành viên.

Hà gửi từng kết quả nhỏ kèm cách kiểm. Thế và Hưng luân phiên rà trước,
đối chiếu chéo các kết quả quan trọng và góp ý để Hà sửa hoặc tích hợp.
Cả ba cùng tham gia rà mã, chứng minh, bài tập và dữ liệu; mỗi thành
viên phải giải thích được bất kỳ đoạn mã hoặc chứng minh nào khi vấn
đáp.

Riêng tạo nhãn DPO, Thế và Hưng đều chấm từng cặp độc lập, không chia
lượt; Hà phân xử sau khi khóa nhãn và đo đồng thuận. Nhóm cần bố trí
lịch và tập huấn trước khi chấm chính thức. Đây là cách tổ chức dự kiến,
chưa phải bằng chứng đã có nhãn hoặc kết quả đồng thuận.

\section{6. Tiến độ và mốc bàn
giao}\label{tiux1ebfn-ux111ux1ed9-vuxe0-mux1ed1c-buxe0n-giao}

{\def\LTcaptype{none} % do not increment counter
\begin{longtable}[]{@{}
  >{\raggedright\arraybackslash}p{(\linewidth - 4\tabcolsep) * \real{0.3333}}
  >{\raggedright\arraybackslash}p{(\linewidth - 4\tabcolsep) * \real{0.3333}}
  >{\raggedright\arraybackslash}p{(\linewidth - 4\tabcolsep) * \real{0.3333}}@{}}
\toprule\noalign{}
\begin{minipage}[b]{\linewidth}\raggedright
Thời điểm
\end{minipage} & \begin{minipage}[b]{\linewidth}\raggedright
Công việc cần hoàn thành
\end{minipage} & \begin{minipage}[b]{\linewidth}\raggedright
Bằng chứng cần lưu
\end{minipage} \\
\midrule\noalign{}
\endhead
\bottomrule\noalign{}
\endlastfoot
Tuần 1, đầu--giữa tuần & Chốt địa phương, nhóm thủ tục, quyền dùng
nguồn; thiết kế schema/rubric/split; khởi tạo mô hình ký tự và DPO
custom & đề cương, data plan, sổ nguồn, mẫu record, quyết định phạm
vi \\
Tuần 1, cuối tuần & Thu thập/chuẩn hóa pilot hợp lệ; kiểm citation trên
mẫu; SFT pilot sớm nếu dữ liệu tối thiểu sẵn sàng & dữ liệu pilot có
provenance, log kiểm tra, checkpoint pilot hoặc ghi rõ blocker \\
Tuần 2, đầu & Hoàn thành hoặc tiến gần ngưỡng dữ liệu D theo tiến độ
thực; dùng checkpoint SFT để tạo ứng viên DPO; chấm độc lập và phân xử &
version dataset, nhãn thô + quyết định phân xử, agreement \\
Tuần 2, theo mốc đề bài & Hoàn thành Phần A, B và bộ dữ liệu Phần D theo
yêu cầu báo cáo tiến độ & mã/phân tích A--C, datasheet, số lượng đã kiểm
chứng (không suy diễn) \\
Tuần 3 & Run cuối theo ngân sách đã chấp thuận, chấm mù/test, báo
cáo/model card/tái lập và chuẩn bị bảo vệ & manifest run, kết quả + CI,
model card, repository/tập lệnh \\
\end{longtable}
}

Các mốc phản ánh quy định đề bài: tuần 1 nộp đề cương và kế hoạch dữ
liệu, tuần 2 báo cáo tiến độ với A/B/dữ liệu D, tuần 3 nộp và bảo vệ.
Tuy nhiên phần cấu trúc chung cũng mô tả nhóm 3 sinh viên thực hiện
trong ``2 tuần'', nên mốc tuần 3 tạo mâu thuẫn lịch trong chính đề bài.
Kế hoạch này không tự coi tuần 3 là gia hạn: nhóm phải xác nhận với
giảng viên mốc áp dụng và báo cáo trung thực phần nào chưa hoàn thành.

\section{7. Rủi ro, giảm thiểu và quyết định cần
chốt}\label{rux1ee7i-ro-giux1ea3m-thiux1ec3u-vuxe0-quyux1ebft-ux111ux1ecbnh-cux1ea7n-chux1ed1t}

{\def\LTcaptype{none} % do not increment counter
\begin{longtable}[]{@{}
  >{\raggedright\arraybackslash}p{(\linewidth - 4\tabcolsep) * \real{0.3333}}
  >{\raggedright\arraybackslash}p{(\linewidth - 4\tabcolsep) * \real{0.3333}}
  >{\raggedright\arraybackslash}p{(\linewidth - 4\tabcolsep) * \real{0.3333}}@{}}
\toprule\noalign{}
\begin{minipage}[b]{\linewidth}\raggedright
Rủi ro
\end{minipage} & \begin{minipage}[b]{\linewidth}\raggedright
Tác động
\end{minipage} & \begin{minipage}[b]{\linewidth}\raggedright
Giảm thiểu/kích hoạt quyết định
\end{minipage} \\
\midrule\noalign{}
\endhead
\bottomrule\noalign{}
\endlastfoot
Phạm vi nguồn quá rộng hoặc văn bản thay đổi & Citation lỗi thời, khối
lượng không kiểm soát & Chốt phạm vi hẹp; lưu version/effective
date/span; không trộn phiên bản không xác định \\
Không đủ quyền/điều khoản thu thập & Không thể dùng nguồn hoặc khảo sát
& Kiểm tra trước ingest; tuân thủ robots/rate limit; dừng nguồn không rõ
quyền; không thu người dùng khi thiếu consent \\
2.000 QA/500 cặp không kịp kiểm chất lượng & Không đạt yêu cầu dữ liệu
hoặc nhãn yếu & Pilot pipeline sớm, kiểm thủ công theo tỷ lệ, tách
validation; báo cáo count thực, không bịa đủ ngưỡng \\
Thiếu hai người chấm độc lập hoặc người khảo sát & Preference/test không
đúng thiết kế & Chốt người chấm và khả năng khảo sát ngay tuần 1; nếu
không có consent/người dùng, ghi rõ không có test người dùng thật thay
vì gán nhãn giả \\
T4/QLoRA hết bộ nhớ hoặc run gián đoạn & Không chạy được SFT/DPO đáng
tin cậy & Pilot SFT rồi DPO với ref riêng; giảm cấu hình sau
validation/pilot có log; lưu checkpoint/manifest để resume khi hạ tầng
cho phép \\
5 run vượt ngân sách & Không đáp ứng báo cáo tái lập & Dự trù 5 seed
ngay từ thiết kế; nếu không đủ, xin điều chỉnh với giảng viên và ghi rõ,
không đổi âm thầm hay chỉ dựa vào đồng thuận nội bộ \\
Rò rỉ split/gần trùng & Test quá lạc quan & Nhóm theo thủ tục/nguồn, lọc
exact/near duplicate và audit citation trước khóa test \\
\end{longtable}
}

Phạm vi Hà Nội/hộ tịch--chứng thực và phương án khảo sát/chấm đã được
đồng ý. Đã có danh mục nguồn sơ bộ, 20 QA nháp nghiên cứu và bộ biểu
mẫu/protocol. Cần giải quyết tiếp: khác biệt thời hạn/phiên bản nguồn,
nguồn địa phương/biểu phí còn thiếu, quyền dùng nội dung, review QA bởi
thành viên, xác nhận vận hành khảo sát, lịch, ngân sách Colab và mô hình
nền sau pilot. Chưa có khảo sát, dữ liệu DPO hoặc benchmark hoàn chỉnh.

\section{8. Tài nguyên và sản phẩm dự
kiến}\label{tuxe0i-nguyuxean-vuxe0-sux1ea3n-phux1ea9m-dux1ef1-kiux1ebfn}

Đề bài khuyến nghị GPU từ 16 GB cho Đề tài 9. Cấu hình T4 nêu ở trên là
thông tin báo cáo từ phiên Colab trước và gần ngưỡng này, nên được xem
là điều kiện cần pilot chứ không phải năng lực đã xác nhận. Tài nguyên
lưu trữ cần chứa bản nguồn hợp lệ, dữ liệu đã ẩn danh,
manifest/checkpoint và log; quyền Drive, dung lượng, giữ phiên và chính
sách resume vẫn chưa chốt.

Sản phẩm cuối theo đề bài dự kiến gồm báo cáo LaTeX 25--40 trang,
repository Git có mô tả môi trường và script tái lập, dữ liệu hoặc
script thu thập kèm datasheet và model card. Đề cương này không tuyên bố
các sản phẩm đó đã tồn tại.

\section{9. Tài liệu căn cứ}\label{tuxe0i-liux1ec7u-cux103n-cux1ee9}

\begin{enumerate}
\def\labelenumi{\arabic{enumi}.}
\tightlist
\item
  TS. Phạm Văn Khánh, \emph{Danh mục 20 bài tập lớn: Học máy cơ bản, mô
  hình ngôn ngữ lớn, học tăng cường và tích hợp}, ngày 12/11/2026: Phần
  I--II (tr. 1--2) và Đề tài 9 (tr. 18--19).
\item
  \href{data-plan.pdf}{Kế hoạch dữ liệu} dẫn tới sổ nguồn, báo cáo 20 QA
  nháp, audit Truyện Kiều và bộ đồng ý/khảo sát/chấm. Những bằng chứng
  kỹ thuật này không thay cho thẩm định quyền/hiệu lực, kiểm thủ công
  hoặc kết quả huấn luyện.
\end{enumerate}
\end{document}
